
% !Mode:: "TeX:UTF-8"
%===Head===
\documentclass[12pt,a4paper,UTF8]{report}
\usepackage{ctex}
\usepackage{CJK}
\usepackage{ulem}
\usepackage{amsmath,amscd,amsbsy,amssymb,latexsym,url,bm,amsthm}
\usepackage{colortbl}
\usepackage{array}
\usepackage{appendix}

\usepackage{epsfig,graphicx,subfigure,wrapfig}
\usepackage{enumitem}
\usepackage{wrapfig}
\usepackage{mathrsfs, euscript}
\usepackage[usenames]{xcolor}
\usepackage{hyperref}
\usepackage{epstopdf}
\usepackage{float}
\usepackage{comment}
\usepackage{listings}
\usepackage{xcolor}
%\usepackage{eso-pic}
\usepackage{graphicx, fancyhdr}
%\usepackage{tocloft}
\usepackage[top=1cm, bottom=1cm, left=2cm, right=2cm]{geometry}
%\usepackage{geometry}
\usepackage{layout}

\usepackage{url}
\makeatletter
\def\UrlAlphabet{%
      \do\a\do\b\do\c\do\d\do\e\do\f\do\g\do\h\do\i\do\j%
      \do\k\do\l\do\m\do\n\do\o\do\p\do\q\do\r\do\s\do\t%
      \do\u\do\v\do\w\do\x\do\y\do\z\do\A\do\B\do\C\do\D%
      \do\E\do\F\do\G\do\H\do\I\do\J\do\K\do\L\do\M\do\N%
      \do\O\do\P\do\Q\do\R\do\S\do\T\do\U\do\V\do\W\do\X%
      \do\Y\do\Z}
\def\UrlDigits{\do\1\do\2\do\3\do\4\do\5\do\6\do\7\do\8\do\9\do\0}
\g@addto@macro{\UrlBreaks}{\UrlOrds}
\g@addto@macro{\UrlBreaks}{\UrlAlphabet}
\g@addto@macro{\UrlBreaks}{\UrlDigits}
\makeatother

\lstset{
    basicstyle = \small,
    numbers=left,
    numberstyle= \tiny,
    keywordstyle= \color{ blue!70},
    commentstyle= \color{red!50!green!50!blue!50},
    frame=shadowbox, % 阴影效果
    rulesepcolor= \color{ red!20!green!20!blue!20} ,
    escapeinside=``, % 英文分号中可写入中文
    xleftmargin=2em,xrightmargin=2em, aboveskip=1em,
    framexleftmargin=2em
}

\newtheorem{theorem}{Theorem}[section]
\newtheorem{lemma}[theorem]{Lemma}
\newtheorem{proposition}[theorem]{Proposition}
\newtheorem{corollary}[theorem]{Corollary}
\newtheorem{exercise}{Exercise}[section]
\newtheorem*{solution}{Solution}

\renewcommand{\thefootnote}{\fnsymbol{footnote}}

\newcommand{\postscript}[2]
 {\setlength{\epsfxsize}{#2\hsize}
  \centerline{\epsfbox{#1}}}

\newsavebox{\mygraphic}
    \sbox{\mygraphic}{%
        \includegraphics[height = 50mm, width = 55mm]{ACM_LOGO_WM.png}
}
\pagestyle{fancy}
\fancyhead[C]{%
    \rhead{\small \leftmark}
    \setlength{\unitlength}{1in}
              \begin{picture}(0, 0)
              \put(1.8, -2){\usebox{\mygraphic}}
              \end{picture}
}
\fancypagestyle{plain}{%
    \rhead{\small \leftmark}
    \fancyhead[C]{\setlength{\unitlength}{1in}
              \begin{picture}(0, 0)
              \put(1.8, -2){\usebox{\mygraphic}}
              \end{picture}}
}

\renewcommand{\baselinestretch}{1.0}

\setlength{\oddsidemargin}{-0.365in}
\setlength{\evensidemargin}{-0.365in}
\setlength{\topmargin}{-0.3in}
\setlength{\headheight}{0in}
\setlength{\headsep}{0in}
\setlength{\textheight}{10.1in}
\setlength{\textwidth}{7in}
\makeatletter \renewenvironment{proof}[1][Proof] {\par\pushQED{\qed}\normalfont\topsep6\p@\@plus6\p@\relax\trivlist\item[\hskip\labelsep\bfseries#1\@addpunct{.}]\ignorespaces}{\popQED\endtrivlist\@endpefalse} \makeatother
\makeatletter
\renewenvironment{solution}[1][Solution] {\par\pushQED{\qed}\normalfont\topsep6\p@\@plus6\p@\relax\trivlist\item[\hskip\labelsep\bfseries#1\@addpunct{.}]\ignorespaces}{\popQED\endtrivlist\@endpefalse} \makeatother

%------------------------------------------------------------------------
%字体设置
\newcommand{\song}{\CJKfamily{zhsong}}
\newcommand{\kai}{\CJKfamily{zhkai}}
\newcommand{\hei}{\CJKfamily{zhhei}}
\newcommand{\fs}{\CJKfamily{zhfs}}

%------------------------------------------------------------------------
%字号设置
\newcommand{\chuhao}{\fontsize{42pt}{\baselineskip}\selectfont}
\newcommand{\xiaochuhao}{\fontsize{36pt}{\baselineskip}\selectfont}
\newcommand{\yihao}{\fontsize{28pt}{\baselineskip}\selectfont}
\newcommand{\xiaoyihao}{\fontsize{24pt}{\baselineskip}\selectfont}
\newcommand{\erhao}{\fontsize{21pt}{\baselineskip}\selectfont}
\newcommand{\xiaoerhao}{\fontsize{17.28pt}{\baselineskip}\selectfont}
\newcommand{\sanhao}{\fontsize{15.75pt}{\baselineskip}\selectfont}
\newcommand{\sihao}{\fontsize{14pt}{\baselineskip}\selectfont}
\newcommand{\xiaosihao}{\fontsize{12pt}{\baselineskip}\selectfont}
\newcommand{\wuhao}{\fontsize{10.5pt}{\baselineskip}\selectfont}
\newcommand{\xiaowuhao}{\fontsize{9pt}{\baselineskip}\selectfont}
\newcommand{\liuhao}{\fontsize{7.875pt}{\baselineskip}\selectfont}
\newcommand{\qihao}{\fontsize{5.25pt}{\baselineskip}\selectfont}

\begin{document}
\begin{titlepage} % Suppresses headers and footers on the title page

    \includegraphics[height = 50mm, width = 55mm]{ACM_LOGO.png}

	\centering % Centre everything on the title page

	\vspace*{5\baselineskip} % White space at the top of the page
	
	\rule{\textwidth}{1.6pt}\vspace*{-\baselineskip}\vspace*{2pt} % Thick horizontal rule
	\rule{\textwidth}{0.4pt} % Thin horizontal rule
	
	\vspace{\baselineskip} % Whitespace above the title
	
	\fs\xiaochuhao{火车站订票系统开发报告} % Title
	
	%\vspace{0.5\baselineskip} % Whitespace below the title
	
	\rule{\textwidth}{0.4pt}\vspace*{-\baselineskip}\vspace{3.2pt} % Thin horizontal rule
	\rule{\textwidth}{1.6pt} % Thick horizontal rule
	
	\vspace{15\baselineskip} % Whitespace after the title block
	
	
    \xiaoerhao
    \fs 组长：\kai\underline{\ \ \ \ \ \ \ \ \ \ \ 郑文鑫\ \ \ \ \ \ \ \ \ \ \ \ }

    \vspace{\baselineskip}
    \fs 组员：\kai\underline{凌皓煜、何俊、赵寒烨}

    \vspace{\baselineskip}
    \fs 日期：\kai\underline{\ \ \ \ 2018年6月15日\ \ \ \ }
\end{titlepage}
%===========================================================================================
\tableofcontents
\newpage
%========================================================================CHAPTER 1===========================================================
\chapter{项目简介}
\section{项目概念}
\par 本项目是上海交通大学2017级ACM班2017-2018学年春季学期数据结构课程的大作业，要求完成类12306的火车票查询与预订系统。其中，后端应支持用户、火车、车票三大部分相关的命令；前端应在Windows和Linux端、网页端、Android端中实现至少一个；前后端使用socket通信。
\par 本次作业后端强制使用C++，且不得使用STL、string、algorithm等工具，要求使用B/B+/B*树一类数据处理形式；前端和前后端通信语言不限。
\par 具体要求请见\color{blue}\url{https://github.com/camelop/ticket_office_cw-DS2018}\color{black}。
\section{项目组成}
\subsection{前端部分}
%TODO
%picture needed
\subsection{后端部分}
\par 后端由解释器、数据库和一些必要的辅助工具组成。
\par 解释器对传入的命令进行解释，并调用相应的数据库接口，再将返回值以一定的格式输出。数据库的主体是B*树，再用多棵B*树分别存储用户、火车等信息。辅助工具包括存储字符串的String、 记录日期时间的Date和Time、传递信息的Vector、Set、Pair等。
\subsection{通信部分}
%TODO
\newpage
%========================================================================CHAPTER 2===========================================================
\chapter{使用手册}
\newpage
%========================================================================CHAPTER 3===========================================================
\chapter{开发日志}
\par 以下列出开发过程中的重要节点，详细记录请访问\color{blue}\url{https://github.com/peterzheng98/2017-2018-B_DS_TrainTicketBooking}\color{black}查看。
\section{前端部分}
%TODO
\section{后端部分}
\setcounter{secnumdepth}{-2}
\subsection{04/14 — \ 05/07}
\par 经过长时间的商讨，敲定了标准接口表。同时初步完成了辅助工具。其中，String、Vector、Pair由郑文鑫编写，Set由赵寒烨编写，凌皓煜进行了测试。
\par 此时，我们错误地认为需要处理中文编码问题，在如何用String正常存储显示中文上耗费了太多时间。
\subsection{05/08 — \ 05/12}
\par 赵寒烨开始动手B*树，第一个能通过编译的版本在5月12日完成。其间修复了String的部分漏洞。
\subsection{05/13 — \ 05/19}
\par 在此期间，后端的主要工作是调试B*树，并于5月26日自以为调试完成。
\subsection{05/21}
\par 郑文鑫开始着手写解释器。
\subsection{05/22 — \ 05/27}
\par 本组唯一能连接远端服务器的电脑损坏，后端开发停滞。
\subsection{05/27 — \ 05/29}
\par 数据库的第一版接口基本完成，但B*树的代码突然在已经通过的测试中崩溃。
\subsection{05/29 — \ 06/05}
\par 赵寒烨修复了B*树的漏洞，但牺牲了大量性能。作为补偿，尝试对B*树的某些部分进行性能优化。
\par 同时完善了第一版接口。
\subsection{06/06}
\par 第一版接口和解释器通过了正确性测试，但性能远远不能满足需求。Test 8在阿里云学生机上运行时，由于占用空间超过2.5G，会导致段异常。
\subsection{06/07}
\par 上午，赵寒烨对第一版接口进行了大幅重构得到第二版接口，性能却仍然太差。Test 8在上述环境测试能得到正确答案，但用时超过600秒。
\par 下午赵寒烨再次重构，并进行两次常数优化，终于把Test 8的运行时间缩短到约34秒，但OJ测试由于不明原因答案错误。
\subsection{06/08}
\par 将所有代码中的部分变量类型改成int32\_t之类的跨平台类型，通过了后端测试。至此后端开发完成。
\newpage
%========================================================================CHAPTER 4===========================================================
\setcounter{secnumdepth}{0}
\chapter{后记}
%TODO
\newpage
%========================================================================APPENDIX===========================================================
\setcounter{secnumdepth}{-2}
\chapter{附录}
%TODO
\end{document}
